\section*{Hoofdstuk \ref{ch:g10:heart-lungs-effort} --- Hart en longen tijdens inspanning}

\begin{solution}{exo:g10:heart-lungs-effort:1}
$0.6 \times 14 = \qty{8.4}{L/min}$; $2.2 \times 35 = \qty{77}{L/min}$.
\end{solution}

\begin{solution}{exo:g10:heart-lungs-effort:2}
Het volume bloed dat één kant van het hart per minuut pompt:
$Q = f \times V_s$. $65 \times 0.075 \approx \qty{4.9}{L/min}$.
\end{solution}

\begin{solution}{exo:g10:heart-lungs-effort:3}
Rechterboezem, rechterkamer, longslagader, longen (kleine bloedsomloop),
longaderen, linkerboezem, linkerkamer, aorta, organen (grote
bloedsomloop), aderen, rechterboezem.
\end{solution}

\begin{solution}{exo:g10:heart-lungs-effort:4}
Ademvolume \qty{0.5}{L}; vitale capaciteit \qty{4.5}{L}; restvolume
\qty{1.2}{L}.
\end{solution}

\begin{solution}{exo:g10:heart-lungs-effort:5}
Ongeveer $220 - 16 = 204$ en $220 - 60 = 160$ slagen per minuut.
\end{solution}

\begin{solution}{exo:g10:heart-lungs-effort:6}
\qty{100}{W}: $120 \times 0.105 \approx \qty{12.6}{L/min}$;
\qty{300}{W}: $198 \times 0.115 \approx \qty{22.8}{L/min}$; een factor
1.8.
\end{solution}

\begin{solution}{exo:g10:heart-lungs-effort:7}
Rust: \qty{2.5}{L/min}; inspanning: $0.04 \times 24 \approx
\qty{1.0}{L/min}$. Het aandeel daalt twaalfvoudig, de absolute stroom
slechts met een factor 2.6, omdat het minuutvolume zelf is gegroeid.
\end{solution}

\begin{solution}{exo:g10:heart-lungs-effort:8}
$18 \times 0.14 = \qty{2.5}{L/min}$: ongeveer drie kwart van de
\qty{3.45}{L/min} van een getrainde leerling --- dicht daarbij voor een
kleiner persoon.
\end{solution}

\begin{solution}{exo:g10:heart-lungs-effort:9}
De hersenen kunnen zelfs geen seconden zonder zuurstof en glucose, dus
wordt hun aanvoer beschermd. De darm en de nieren kunnen de vertering en
de filtering een uur stilleggen zonder schade, en hun aandeel wordt aan
de spieren uitgeleend.
\end{solution}

\begin{solution}{exo:g10:heart-lungs-effort:10}
$V_s = 5000/42 \approx \qty{120}{mL}$, tegenover ongeveer \qty{70}{mL}:
haar vergrote hart levert het minuutvolume in rust met veel minder
slagen.
\end{solution}

\begin{solution}{exo:g10:heart-lungs-effort:11}
Het bevel van de hersenen om te bewegen, en hun vooruitlopen op de
inspanning, worden naar de hartcentra en de bijnieren gestuurd voordat de
spieren in actie komen. Een loper die alleen op het meten van
koolstofdioxide vertrouwde, zou met een minuutvolume in rust starten, de
eerste minuut op \omterm{def:g10:cell-metabolism:fermentation}{gisting} lopen, en door melkzuur gedwongen worden te
vertragen voordat de aanvoer bij was.
\end{solution}

\begin{solution}{exo:g10:heart-lungs-effort:12}
$195 \times 0.110 \times 0.150 \approx \qty{3.2}{L/min}$ en
$195 \times 0.160 \times 0.150 \approx \qty{4.7}{L/min}$. Het slagvolume
is wat training vergroot: de tweede proefpersoon is de duurgetrainde.
\end{solution}

\begin{solution}{exo:g10:heart-lungs-effort:13}
Ingeademde zuurstof: $0.21 \times 100 = \qty{21}{L/min}$; gebruikt:
\qty{3.5}{L/min}, dus 17\%. De rest wordt weer uitgeademd --- uitgeademde
lucht bevat nog ongeveer 16\% zuurstof, en daarom werkt
mond-op-mondbeademing.
\end{solution}

\begin{solution}{exo:g10:heart-lungs-effort:14}
Het minuutvolume komt hoogstens op $110 \times 0.12 \approx \qty{13}{L/min}$
in plaats van 24: de $\dot V\!\mathrm{O_2}$max daalt met
bijna de helft, tot ongeveer \qty{2}{L/min}. De patiënt kan rustig lopen
en fietsen, maar geen zware inspanning volhouden.
\end{solution}

\begin{solution}{exo:g10:heart-lungs-effort:15}
De aanvoer per liter daalt met een kwart, de onttrekking blijft
verhoudingsgewijs gelijk: $\dot V\!\mathrm{O_2}$max daalt met 25\%. In de
weken daarna maakt het lichaam meer rode bloedcellen, waardoor de
zuurstof die per liter wordt gedragen weer richting \qty{200}{mL} gaat.
\end{solution}

\begin{solution}{pb:g10:heart-lungs-effort:1}
\textbf{1.} $62 \times 0.072 \approx 4.5$; $105 \times 0.100 = 10.5$;
$140 \times 0.116 \approx 16.2$; $175 \times 0.120 = 21.0$;
$198 \times 0.120 \approx \qty{23.8}{L/min}$.

\textbf{2.} $23.8/4.5 \approx 5.3$: de frequentie draagt $198/62
\approx 3.2$ bij, het slagvolume $120/72 \approx 1.7$, en $3.2 \times
1.7 \approx 5.3$.

\textbf{3.} Boven ongeveer \qty{225}{W} blijft het slagvolume op
\qty{120}{mL}; alleen de hartfrequentie verhoogt het minuutvolume nog.

\textbf{4.} $220 - 17 = 203$: zijn 198 valt binnen de nauwkeurigheid van de regel.

\textbf{5.} $5/4.5 \approx \qty{1.1}{min}$, ongeveer \qty{67}{s}; bij
\qty{300}{W}, $5/23.8 \approx \qty{0.21}{min}$, ongeveer \qty{13}{s}.

\textbf{6.} $0.5 \times 13 = 6.5$; $1.2 \times 18 = 21.6$; $1.8 \times
24 = 43.2$; $2.3 \times 32 = 73.6$; $2.6 \times 44 \approx
\qty{114}{L/min}$.

\textbf{7.} $114/6.5 \approx 18$: het ademvolume steeg 5.2 keer, de
frequentie 3.4 keer; het volume draagt meer bij.

\textbf{8.} $2.6/5.0 = 52\%$ van de vitale capaciteit per ademteug.

\textbf{9.} Nee: de \omterm{def:g10:heart-lungs-effort:ventilation}{ventilatie} heeft een grote reserve en het bloed dat
de longen verlaat blijft verzadigd. De grens is het slagvolume van het
hart, dat bepaalt hoeveel zuurstof elke slag aanvoert.

\textbf{10.} $0.7 \times 6.5 \approx \qty{4.6}{L/min}$ en $0.7 \times
114 \approx \qty{80}{L/min}$.

\textbf{11.} $4.5 \times 0.052 \approx 0.23$; $10.5 \times 0.090
\approx 0.95$; $16.2 \times 0.120 \approx 1.94$; $21.0 \times 0.145
\approx 3.05$; $23.8 \times 0.155 \approx \qty{3.7}{L/min}$.

\textbf{12.} $3700/68 \approx \qty{54}{mL/min}$ per kilogram. De
frequentie heeft haar maximum bereikt en de toename van 225 naar
\qty{300}{W} is klein: dit is zijn $\dot V\!\mathrm{O_2}$max, die van een
fitte, getrainde leerling.

\textbf{13.} $3.7/0.23 \approx 16$: de factor 5.3 van het hart
vermenigvuldigd met de $155/52 \approx 3.0$ van de onttrekking.

\textbf{14.} $3.7 \times 20 = \qty{74}{kJ}$ per minuut, ongeveer
\qty{1230}{W}; in rust $0.23 \times 20/60 \approx \qty{77}{W}$; boven de
rust ongeveer \qty{1150}{W}, dus is het rendement $300/1150 \approx
26\%$.

\textbf{15.} Ruwweg een rechte lijn van \num{0.23} naar
\qty{3.05}{L/min} tussen 0 en \qty{225}{W}, daarna buigend: de laatste
\qty{75}{W} voegen slechts \qty{0.65}{L/min} toe. Het plafond ligt rond
\qty{3.7}{L/min}.

\textbf{16.} Vooruitlopen: de verwachting van de inspanning bij de
hersenen wordt aan de hartcentra en de bijnieren doorgegeven, en verhoogt
de frequentie nog vóór enige beweging.

\textbf{17.} De \omterm{def:g10:heart-lungs-effort:ventilation}{longblaasjes} voorzien het bloed zelfs bij maximale
\omterm{def:g10:heart-lungs-effort:ventilation}{ventilatie} en maximaal minuutvolume volledig van zuurstof; de longen zijn
niet de grens. De grens is de hoeveelheid bloed die het hart per minuut
kan versturen.

\textbf{18.} $20/23.8 \approx 84\%$. De darm en de nieren, die de helft
van het minuutvolume in rust namen, ontvangen nu ongeveer \qty{1}{L/min};
hun slagaders zijn vernauwd en de stroom is naar de spieren geleid.

\textbf{19.} $198 \times 0.140 \times 0.155 \approx \qty{4.3}{L/min}$,
ongeveer \qty{63}{mL/min} per kilogram.

\textbf{20.} Het \omterm{def:g10:heart-lungs-effort:output}{hartminuutvolume} vermenigvuldigd met 5.3, de
zuurstofonttrekking met 3, de \omterm{def:g10:heart-lungs-effort:ventilation}{ventilatie} met 18; de training veranderde
het slagvolume van het hart, en daarmee het plafond.
\end{solution}
